% Auxiliary input to the maps-owned OS45 fundamental-class calculation.
% The full boundary comparison identifying the support with global u is required.
\subsubsection*{The remaining finite-coefficient evaluations for OS45}
Let $u\in H^2(X;\mathbb Z/8)$ be the actual primitive torsion
coordinate constructed by the complete boundary comparison, with
cusp support
\[
 \Psi(r,s)=4\left\lfloor\frac r8\right\rfloor
       +3(r\bmod8)\left\lfloor\frac s2\right\rfloor
                           \quad\pmod8.                  \tag{D.OL1}
\]
Here $r\bmod8$ is the representative from zero to seven. This
subsection uses that comparison on the whole cusp boundary;
matching only the regular-fibre bilinear form would not suffice.
The actual affine changes are
\[
 \Delta_{(8,-6)}\Psi=-r+4,\qquad
 \Delta_{(24,-16)}\Psi=4\quad\pmod8.                     \tag{D.OL2}
\]
In particular both have zero slope in the scalar $s$ direction.
Let $y$ be the primitive class in (D.PC9).
We prove the three full evaluations
\[
 \langle u^3,[X]\rangle=0\pmod8,\qquad
 \langle\mathcal P_2(u)y,[X]\rangle=0\pmod{16},\qquad
 \langle uy^2,[X]\rangle=0\pmod8.                         \tag{D.OL3}
\]
The second formula includes the order-two part which cannot be
recovered by halving the sum of the component evaluations.

The scalar circle action is semifree. Its fixed locus consists of
the sixteen vertical toric curves $C_{r,s}$ joining vertices
$(r,s)$ and $(r,s+1)$, with $(r,s)$ indexed as in (D.PC6).
Each has normal bundle $\mathcal O(-1)\oplus\mathcal O(-1)$
with weights $1,-1$. In every other triangular toric chart the
nonzero weights are also primitive. At each of the four zero-free
$I_1$ caps the primitive angular collapsing vector has first
coordinate one in its elliptic/scalar plane; the scalar circle
therefore projects primitively to the transverse torus and is
free. The ordinary torus cap is free as well. Thus there are no
unrecorded finite-stabilizer strata in this localization.

Since $X$ is simply connected, both $u$ and $y$ have circle-equivariant
lifts. Indeed the only outgoing differentials from ordinary degree
two in the Borel spectral sequence have targets involving
$H^1(X)$ or $H^3(BS^1)$, both zero. Let $x$ be the positive
generator of $H^2(C_{r,s};\mathbb Z)$ and let $t$ be the generator
of $H^2(BS^1;\mathbb Z)$. Write
\[
 \widehat u|_{C_{r,s}}=d_{r,s}x+\lambda_{r,s}t,
 \qquad
 \widehat y|_{C_{r,s}}=\mu_{r,s}t.                         \tag{D.OL4}
\]
There is no $x$ term in the second expression: (D.PC14) and
$f\cdot C_{r,s}=0$ give $y(C_{r,s})=0$.

For clarity the ordinary degrees and the fibre characters are
computed separately. The degree of the support on a vertical wall is
\[
 d_{r,s}=\Psi(r,s)+\Psi(r,s+1)
              -\Psi(r+1,s)-\Psi(r-1,s+1)\pmod8 .         \tag{D.OL5}
\]
Choosing representatives from zero to seven gives
\[
\begin{array}{c|cc}
 r&d_{r,0}&d_{r,1}\\\hline
 0,7&4&7\\
 1,\ldots,6&0&3.
\end{array}                                                \tag{D.OL6}
\]
The scalar fibre character supplied by the support is its vertical
slope, up to a single common sign. Thus one may take
\[
                 \lambda_{r,0}=0,\qquad
                 \lambda_{r,1}=3r .                      \tag{D.OL7}
\]
This is periodic under the actual period lattice because of
(D.OL2). To justify the common character, rather than making a
separate choice on each normalized component, compare along the
actual toric orbit curves. Their degrees determine the differences
of the characters at their fixed endpoints. Horizontal and diagonal
orbit curves connect the sixteen fixed curves in a connected graph:
they join $(r,s)$ to $(r+1,s-1)$ and to $(r-1,s)$ respectively.
The support realizes all these differences. Consequently any global
equivariant lift has the characters in (D.OL7), with a common sign
convention and one common additive character $L\in\mathbb Z/8$.
No global integral line lifting $u$ is assumed.

The integral component divisors have their canonical equivariant
linearizations. Along $C_{r,s}$ the two adjacent component normal
weights are $-1,+1$ in the lower/upper vertex order. Thus the weight
of the odd-component sum in (D.PC9) is $+1$ for $s=0$ and $-1$
for $s=1$. The difference
$2\widehat y-\sum_r\widehat x_{r,1}$ is a global multiple $ct$.
The existence of integral fixed-point weights forces $c$ odd;
changing the lift of $y$ changes $c$ by an even integer. Normalize
$c=-1$. Then
\[
                       \mu_{r,0}=0,\qquad
                       \mu_{r,1}=-1.                     \tag{D.OL8}
\]
Every other lift adds one common integer $V$.

Each restriction of $\widehat u$ in (D.OL4) lifts integrally on
$C_{r,s}\times BS^1$. Thomas's cochain formula therefore gives
$\mathcal P_2(\widehat u)=(d x+\lambda t)^2$ with the full
modulus sixteen \cite[equations (2.1), (3.4), pp.267--268]{ffThomas56}.
The ordinary cube uses modulus eight. The equivariant normal Euler
class is
\[
                           (-x+t)(-x-t)=-t^2.
\]
Localization is valid over these coefficient rings because the
complement of the fixed locus is free and all normal weights are
units. Taking the coefficient of $x$ after division by $-t^2$ gives
\[
 \begin{aligned}
 \langle u^3,[X]\rangle&=-3\sum_{r,s}d_{r,s}\lambda_{r,s}^2
                                                        &&\pmod8,\\
 \langle\mathcal P_2(u)y,[X]\rangle
       &=-2\sum_{r,s}d_{r,s}\lambda_{r,s}\mu_{r,s}
                                                        &&\pmod{16},\\
 \langle uy^2,[X]\rangle&=-\sum_{r,s}d_{r,s}\mu_{r,s}^2
                                                        &&\pmod8.
 \end{aligned}                                             \tag{D.OL9}
\]
The degree-zero base ring injects into its localization, so these
are identities before inverting $t$ as well.

The representatives (D.OL6)--(D.OL8) satisfy the exact integer sums
\[
 \sum d_{r,s}=40,\qquad \sum d_{r,s}\lambda_{r,s}=336,
 \qquad \sum d_{r,s}\lambda_{r,s}^2=5544,
 \qquad \sum_r d_{r,1}=32.                               \tag{D.OL10}
\]
The three values in (D.OL9) are consequently $-16632$, $672$,
and $-32$, proving (D.OL3) with their full moduli.
For completeness, under the replacements
$\lambda\mapsto L+\epsilon\lambda$ and $\mu\mapsto V+\mu$,
where $\epsilon=\pm1$, the exact values are
\[
 \begin{aligned}
 -120L^2-2016\epsilon L-16632,\\
 -80LV+64L-672\epsilon V+672\epsilon,\\
 -40V^2+64V-32.
 \end{aligned}                                             \tag{D.OL11}
\]
Every coefficient is divisible by eight, sixteen, and eight,
respectively. Thus the results are independent of all global
linearization choices, not only of the displayed normalization.

The same support gives all component mixed evaluations. For a
vertex $v$, let $R_j$ be the six cyclic rays of D.2.47 and put
\[
 d_j(v)=\Psi(v)+\Psi(v+R_j)
                  -\Psi(v+R_{j-1})-\Psi(v+R_{j+1}).
\]
An integral representative of its component restriction is
\[
 \eta_v=(d_0+d_1+d_2)H-d_0E_1-d_2E_2-d_4E_3.
                                                               \tag{D.OL12}
\]
Changing the representative by eight times an integral class
changes neither its square modulo sixteen nor any of its needed
pairings modulo eight. In the primitive basis of (D.PC10),
\[
 \langle\mathcal P_2(u)x_i,[X]\rangle=7\pmod{16}
                     \quad(0\le i<14),
\]
while the two omitted components $i=14,15$ have value fifteen.
Also $\langle ux_i y,[X]\rangle$ equals four for $i=0,1,14,15$
and zero for all the other component labels. The remaining
$\langle ux_i x_j,[X]\rangle$ are exactly
$\eta_i\cdot\nu_i^*D_j\pmod8$, with the explicit Cartier
restriction in D.2.47. Together with (D.OL3) these determine the
finite-coordinate evaluations once the support comparison is
established.

The exact program
\path{code/fourfolds_OS45_localization.py}
retains every wall value and every $\eta_v$, verifies the actual
quasiperiods on all sixteen residues, and proves the six
all-character polynomial divisibility identities (D.OL11).
It explicitly records the required geometric support comparison
as a separate hypothesis.
